\begin{lemma}[Scheide defective-color setup]
Let a reduced Scheide configuration have at least one defective color.
Then it admits defective-color setup data: a common
$\alpha_0\in\overline\varphi(y_0)\setminus\Gamma$ and, for every defective
$\delta$, the two full incidence paths, last-tree-vertex indices, named
boundary edge, outside endpoint, and missing block color specified in
\noderef{ScheideDefectiveSetup}.
\end{lemma}
\begin{proof}
Both ends of the uncolored edge $e_0$ are among the first two base vertices,
by the position and attachment fields of
\noderef{ScheideReducedConfiguration}.  The missing-count clause
\noderef{ScheideGlobalMaximumMissingCounts} gives at least two colors
missing at $y_0$.  By the exit-color bijection in
\noderef{ScheideReducedConfiguration}, every block color is missing at a
unique exit; elementariness from \noderef{ScheideGlobalMaximumOrder}
therefore implies that $y_0$ misses at most one block color.  Choose a
second missing color $\alpha_0$ outside $\Gamma$.

Fix a defective color $\delta$.  Closedness from
\noderef{ScheideGlobalMaximumOrder}, together with the existence of a
$\delta$ boundary edge, says that $\delta$ is present at every tree
vertex.  The vertex-node $y_0$ has no $\alpha_0$ incidence and one
$\delta$ incidence, so it is an end of its finite two-color component.
Every edge-node has its two distinct endpoint neighbors, and properness
gives vertex degree at most two.  Hence the spanning-path argument of
\noderef{TwoColorComponentSpanningPath} gives a simple incidence path
whose support is the full component.  Its other end is outside $V(T)$:
no other tree vertex can be an end, since $\delta$ is present everywhere
and elementariness makes $\alpha_0$ present away from $y_0$.

Orient this path from $y_0$.  Finiteness gives its last vertex-node in
$V(T)$, call it $v^0_\delta$, followed by an edge-node
$f^0_\delta$ and an outside vertex-node $u^0_\delta$.  No
$\alpha_0$-edge leaves the closed tree, by
\noderef{ScheideGlobalMaximumBoundaryFacts}; consequently
$\varphi(f^0_\delta)=\delta$.  The last-tree-vertex property makes
$v^0_\delta$ a Scheide exit.  The reduced exit-color bijection supplies
the unique index $\gamma(\delta)$ whose block color is missing there.

Finally consider the incidence component in colors
$\gamma(\delta),\delta$ through $v^0_\delta$.  At its root the first color
is missing and the second is present.  Properness again gives a finite path
component.  It has an outside vertex because a defective $\delta$-edge
leaves the closed tree and the alternating component cannot end at a second
tree vertex: that would duplicate a missing color in the elementary set.
Choose its full simple spanning path oriented from $v^0_\delta$.
Performing these finite choices for every defective color supplies all
fields, path supports, and identity equations in
\noderef{ScheideDefectiveSetup}.
\end{proof}
